% !TEX program = tectonic
%
% NeurIPS-format companion paper: call-type probing and clustering study.
%
% This is the ML-community follow-up to the JASA acoustics paper
% (../manuscript.tex), which established a binary call/non-call detector.
% This paper targets funding-proposal Objective 2 (classifying calls into
% morphological sub-types, including unsupervised subdivision of the
% "complex call" catch-all category) as an open representation-learning /
% weak-supervision problem, rather than assuming it is solved.
%
% Compile with:  tectonic ml_manuscript.tex
% (or: pdflatex ml_manuscript && bibtex ml_manuscript && pdflatex x2)
%
% To submit anonymously to a venue, remove the "preprint" option below --
% the style file automatically anonymizes authors and adds line numbers
% in that mode. Use [final] for the camera-ready version.

\documentclass{article}
\usepackage[preprint,nonatbib]{neurips_2023}
\usepackage[numbers]{natbib}

\usepackage[utf8]{inputenc}
\usepackage[T1]{fontenc}
\usepackage{hyperref}
\usepackage{url}
\usepackage{booktabs}
\usepackage{amsfonts}
\usepackage{microtype}
\usepackage{graphicx}
\usepackage{amsmath}

\title{Do Off-the-Shelf Detector Embeddings Reveal Bowhead Whale Call
Sub-Types? A Probing and Clustering Study on Weakly-Labeled Passive
Acoustic Data}

\author{%
  Oc\'eane Boulais \\
  % Affiliation \\
  % \texttt{email@domain} \\
  \And
  Aaron M.\ Thode \\
  Scripps Institution of Oceanography, University of California San Diego \\
}

\begin{document}

\maketitle

\begin{abstract}
Manual bioacoustic taxonomies frequently include a catch-all category for
sounds that resist clean classification; for migrating bowhead whales
(\emph{Balaena mysticetus}), roughly a third of manually annotated calls
across seven morphological "Type" labels fall into categories described by
human analysts as ambiguous or context-dependent. A companion acoustics
paper establishes a convolutional detector that separates bowhead calls
from other transient sounds with high accuracy (ROC-AUC $\approx$ 0.96 on
a fully independent, 199,825-sample generalization test), but that
detector is explicitly binary and does not model call sub-type structure.
Here we ask two narrower questions about the representations available
from that detector's development: (1) how much information about the
seven manually annotated call sub-types is linearly decodable from
embeddings trained only for reconstruction (a convolutional autoencoder)
or for binary call/non-call detection (from-scratch and
autoencoder-warm-started classifiers); and (2) whether unsupervised
clustering of these same embeddings recovers the manual sub-type taxonomy
without using any sub-type labels. We find that a linear probe recovers
moderate sub-type signal (best macro one-vs-rest ROC-AUC 0.663, top-1
accuracy 36.5\% against a 14.3\% random baseline over seven classes), and
that embeddings retaining more reconstructive detail (the autoencoder and
warm-started classifier) are noticeably more informative for sub-type
identity than embeddings optimized purely for binary detection (scratch
classifier: macro ROC-AUC 0.577). In contrast, standard unsupervised
clustering (k-means, BIC-selected Gaussian mixtures) fails to recover the
manual taxonomy in any embedding space we tested (adjusted Rand index
$\leq$ 0.055 in every configuration), indicating that whatever call-type
signal exists in these embeddings is not axis-aligned with simple cluster
geometry. We argue this dissociation -- weak-to-moderate linearly
decodable structure alongside a near-total failure of unsupervised
clustering -- is itself the useful empirical finding: it shows that
representations optimized for reconstruction or binary detection are an
insufficient basis for the unsupervised call-type subdivision envisioned
in prior work, and motivates representation-learning objectives (e.g.,
contrastive or metric-learning pretraining) explicitly targeted at
sub-type structure, which we outline as the next step in this research
program.
\end{abstract}

\section{Introduction}
\label{sec:intro}

Between 2007 and 2014, 35 passive acoustic recorders (Directional Autonomous
Seafloor Acoustic Recorders, DASARs) were deployed along a 280\,km stretch of
the Alaskan North Slope to monitor the fall migration of the
Bering--Chukchi--Beaufort bowhead whale population. Manual analysts
annotated millions of detected calls with a time, frequency range, and one
of seven morphological "Type" labels. A long-standing observation in this
literature, reflected directly in the manual annotation handbook used to
build this dataset, is that a substantial fraction of calls resist clean
categorization into simple frequency-modulated archetypes (up-calls,
down-calls, undulations) and are instead placed into broader,
harder-to-define categories describing mixed frequency/amplitude
modulation with broadband energy. Whether this catch-all grouping reflects
a single underlying vocalization type or several distinct call types that
happen to be difficult for a human analyst to separate by eye has been an
open question for at least two decades, and is precisely the question a
purely supervised classifier -- which can only reproduce the labels it is
given -- cannot answer on its own.

A companion acoustics paper describes a two-stage detector for this
dataset: a convolutional autoencoder that compresses single-channel SNR
spectrograms into a 32-dimensional latent space, and a supervised
classifier (trained either from random initialization, or warm-started
from the autoencoder's encoder weights) that distinguishes bowhead calls
from other transient sounds, including seismic airgun pulses, with
ROC-AUC $\approx$ 0.96 on a fully independent held-out evaluation set. That
detector is deliberately binary, and treats all seven manually annotated
call sub-types identically during training and evaluation. This raises a
natural follow-up question that the detector itself was not designed to
answer: \emph{do the representations produced along the way to building
that binary detector already carry usable information about call
sub-type, and if so, is that information organized in a way that
unsupervised clustering could exploit to automatically discover or refine
sub-type structure without additional manual annotation?}

We treat this as an empirical probing question rather than assume an
answer. Concretely, we study three 32-dimensional embedding spaces that
already existed as artifacts of building the binary detector -- a purely
reconstructive autoencoder embedding, a from-scratch binary-detection
embedding, and an autoencoder-warm-started binary-detection embedding --
and ask (RQ1) how much of the manually annotated seven-way call-type
structure is linearly decodable from each, and (RQ2) whether standard
unsupervised clustering methods (k-means, Gaussian mixtures with
BIC-selected component count) recover that same seven-way structure
without being given any sub-type labels. Our contribution is not a new
architecture, but a carefully controlled empirical answer, using a
leakage-free evaluation protocol inherited directly from the companion
detector paper, to a question that is foundational to any future attempt
at unsupervised or weakly-supervised bowhead call-type discovery: representations trained for
reconstruction or binary detection are demonstrably \emph{not}
sufficient, on their own, for that goal, which motivates the
representation-learning objectives we outline as future work in
Section~\ref{sec:discussion}.

\section{Related work}
\label{sec:related}

\textbf{Bioacoustic deep learning.} Convolutional and, increasingly,
transformer-based architectures have substantially improved automated
detection of marine mammal and avian vocalizations over classical
spectrogram-correlation and hand-engineered feature approaches
\citep{shiu2020deep, allen2021convolutional}. Prior bowhead-specific work
used hand-crafted features and a shallow three-layer network to flag calls
without attempting sub-type classification \citep{thode2012automated}. Our
companion detector paper replaces this with a modern convolutional
detector, but likewise stops short of sub-type classification, which is
the gap this paper addresses.

\textbf{Representation learning for bioacoustics.} Convolutional
autoencoders and related unsupervised encoders have been used to compress
bioacoustic spectrograms into low-dimensional embeddings for downstream
clustering and visualization \citep{ozanich2021deep, hinton2006reducing}.
Most such studies, however, evaluate reconstruction quality or
binary/coarse detection performance rather than directly asking whether
the resulting embedding geometry supports \emph{unsupervised} recovery of
a pre-existing fine-grained taxonomy, which is the specific question we
address here.

\textbf{Weakly-labeled and open-set classification.} Because the "complex
call" category in this dataset was defined by human analysts as a
catch-all for signals that did not fit cleaner categories, it is best
understood as a weak or noisy label rather than a well-defined class. This
connects our study to the broader literature on learning under label
noise and open-set/novel-class discovery, where a central question is
whether representations trained for a related but distinct task (here,
binary detection) transfer usefully to separate finer-grained or entirely
novel classes not explicitly supervised during training.

\section{Data and problem setup}
\label{sec:data}

We reuse, without modification, the leakage-controlled 100{,}000-sample
dataset and grouped date$\times$site train/test split described in the
companion detector paper: 50{,}000 manually verified bowhead calls (Types
1--7) and 50{,}000 auto-detected non-call transients (Type 0), drawn from
DASAR sites 3 and 5 and DASARs A, D, and G. Every image is a 121$\times$104
single-channel SNR spectrogram spanning 3~seconds and 25--500~Hz, centered
on the time of peak local power spectral density (see companion paper for
full preprocessing details). Splitting is grouped by the
date$\times$site key used to record each detection, so that all images
from a given site on a given day fall entirely within one partition --
the same leakage-safety invariant used throughout the companion paper --
yielding 91{,}920 training images (43 groups) and 8{,}080 held-out test
images (8 groups), with zero group overlap between partitions.

For the present study we restrict attention to the 50{,}000 manually
verified calls, which carry one of seven Type labels (Table~\ref{tab:counts}).
We treat these Type labels exactly as delivered by the original
manual-analysis pipeline: as opaque integer categories, with no semantic
mapping (e.g. "up-call," "complex call") reconstructed or assumed, since no
such mapping exists in the dataset artifacts available to us. This is a
deliberate choice: any claim about which Types correspond to "simple" or
"complex" calls in the sense used by the original annotation handbook would
require information we do not have, so we evaluate all seven Types
identically and let the data speak for itself. Restricting to the grouped
split above yields 47{,}400 training and 2{,}600 test calls with Type
labels (Table~\ref{tab:counts}); Type 6 is the rarest (1{,}553 examples
total) and Type 1 the most common (14{,}693 examples total), a roughly
9.5:1 imbalance ratio that we account for using class-balanced probe
training (Section~\ref{sec:methods}).

\begin{table}[t]
  \centering
  \caption{Manually annotated call-type counts in the grouped train/test
  split used for the probing and clustering study (identical split
  convention as the companion detector paper).}
  \label{tab:counts}
  \begin{tabular}{lccccccc|c}
    \toprule
     & Type 1 & Type 2 & Type 3 & Type 4 & Type 5 & Type 6 & Type 7 & Total \\
    \midrule
    Train & 13{,}951 & 10{,}049 & 7{,}576 & 4{,}548 & 3{,}292 & 1{,}456 & 6{,}528 & 47{,}400 \\
    Test  &    742 &    701 &   392 &   215 &   165 &    97 &   288 &  2{,}600 \\
    \bottomrule
  \end{tabular}
\end{table}

\section{Methods}
\label{sec:methods}

\textbf{Embedding spaces.} We evaluate three 32-dimensional embedding
spaces, all already produced during the construction of the companion
binary detector, with no additional training performed for this study:
\begin{itemize}
  \item \textbf{Autoencoder (AE)}: the encoder trunk of the convolutional
  autoencoder, trained purely to minimize pixel-wise reconstruction error,
  with no access to any call/non-call or Type label during training.
  \item \textbf{CNN (scratch)}: the encoder trunk of the binary
  call/non-call classifier, trained end-to-end from random initialization
  using only the binary label.
  \item \textbf{CNN (warm-start)}: the same binary classifier architecture,
  but with the encoder trunk initialized from the trained autoencoder and
  then fine-tuned end-to-end on the binary label.
\end{itemize}
All three share an identical architecture (three Conv+BatchNorm+ReLU+MaxPool
blocks, FC-64$\to$FC-32 bottleneck), differing only in what objective (if
any) supervised their training, isolating the effect of training objective
on call-type informativeness while holding architecture and input data
fixed.

\textbf{Linear probing (RQ1).} For each embedding space, we extract
32-dimensional embeddings for every manually-verified call in the grouped
train/test split (Section~\ref{sec:data}) and fit a
class-balanced, multinomial logistic-regression probe (StandardScaler
$\to$ balanced multinomial logistic regression, $L_2$ regularization
$C=1$) on the training embeddings and Type labels, following standard
frozen-embedding probing methodology. We report one-vs-rest ROC-AUC
per Type, macro-averaged ROC-AUC across the seven types, and top-1
accuracy, all computed on the held-out grouped test partition. A
macro ROC-AUC of 0.5 and a top-1 accuracy of $1/7 \approx 14.3\%$
correspond to chance performance under this seven-way task.

\textbf{Unsupervised clustering (RQ2).} For each embedding space, we
standard-scale the pooled train+test call embeddings (50{,}000 points) and
apply two clustering algorithms: (i) k-means with $k=7$, matching the
number of manually annotated Types, and (ii) a Gaussian mixture model
with the number of components selected by Bayesian Information Criterion
(BIC) over $k \in \{2,4,6,7,8,10,12,16\}$, allowing the data to select a
different granularity than the manual taxonomy if warranted. We score
cluster quality using the silhouette coefficient (an unsupervised,
label-free measure of cluster separation) and, critically, score
\emph{agreement with the manual Type labels} using the adjusted Rand
index (ARI) and normalized mutual information (NMI) -- two standard
clustering-agreement metrics that are 0 in expectation for a random
partition and 1 for a perfect match to the reference labels.

\section{Results}
\label{sec:results}

\subsection{RQ1: call-type information is linearly decodable, but modest}

Table~\ref{tab:probe} reports the linear-probe results for all three
embedding spaces. All three exceed the chance baseline (macro ROC-AUC
0.5, top-1 accuracy 14.3\%) by a wide margin, confirming that
call-sub-type information is present and linearly accessible in these
embeddings despite none of them being trained with sub-type supervision.
However, the gap between embedding spaces is substantial and
systematic: the autoencoder and warm-started embeddings (macro ROC-AUC
0.659 and 0.663) are both markedly more informative for call sub-type
than the from-scratch binary-detection embedding (macro ROC-AUC 0.577).
This pattern holds for six of seven individual Types
(Table~\ref{tab:probe}), and is consistent with the interpretation that
an embedding optimized purely to separate "call" from "non-call" discards
within-call visual detail that is not needed for that binary decision, but
that a reconstructive or partially-reconstructive objective retains.
Types 1, 2, and 7 are the most separable across all three embeddings
(per-type ROC-AUC 0.72--0.79 for AE/warm-start); Types 3, 5, and 6 are
consistently the hardest to separate (ROC-AUC 0.52--0.60), independent of
which embedding is used.

\begin{table}[t]
  \centering
  \caption{Seven-way call-type linear-probe results (macro one-vs-rest
  ROC-AUC, top-1 accuracy, and per-Type ROC-AUC) on the held-out grouped
  test split (chance: macro ROC-AUC 0.5, top-1 accuracy 14.3\%).}
  \label{tab:probe}
  \small
  \setlength{\tabcolsep}{4pt}
  \begin{tabular}{lccccccccc}
    \toprule
    Embedding & Macro & Top-1 & \multicolumn{7}{c}{Per-type ROC-AUC} \\
    & ROC-AUC & Acc.\ & T1 & T2 & T3 & T4 & T5 & T6 & T7 \\
    \midrule
    Autoencoder    & 0.659 & 36.1\% & 0.735 & 0.722 & 0.564 & 0.619 & 0.596 & 0.584 & 0.794 \\
    CNN (scratch)      & 0.577 & 15.9\% & 0.606 & 0.573 & 0.542 & 0.570 & 0.520 & 0.583 & 0.645 \\
    CNN (warm-start)   & \textbf{0.663} & \textbf{36.5\%} & 0.744 & 0.754 & 0.580 & 0.623 & 0.579 & 0.604 & 0.754 \\
    \bottomrule
  \end{tabular}
\end{table}

\subsection{RQ2: unsupervised clustering fails to recover the manual taxonomy}

Table~\ref{tab:cluster} reports clustering results for all three
embeddings. In every configuration -- both algorithms, all three embedding
spaces -- agreement with the manual seven-way Type labels is close to
zero (ARI $\leq 0.055$, NMI $\leq 0.055$), far below what would indicate
even a loose correspondence between discovered clusters and the manual
taxonomy. This holds despite k-means finding geometrically well-separated
clusters in some embeddings by its own internal criterion (silhouette
0.306 for the scratch-CNN embedding) -- meaning the clusters that do
exist in these embedding spaces are organized around structure other than
manually annotated call type (plausibly SNR, site, DASAR channel, or
signal duration, none of which were controlled for in this analysis). The
BIC criterion for the Gaussian mixture model selected a substantially
larger number of components (12--16) than the seven manually annotated
types in every embedding, and even at that finer granularity, agreement
with the manual Types did not improve materially.

\begin{table}[t]
  \centering
  \caption{Unsupervised clustering agreement with the manual seven-Type
  taxonomy, pooled across the full 50{,}000 manually-verified calls.
  Silhouette is a label-free internal cluster-quality score; ARI and NMI
  measure agreement with the manual Type labels (0 = chance, 1 = perfect).}
  \label{tab:cluster}
  \begin{tabular}{llcccc}
    \toprule
    Embedding & Method & $k$ & Silhouette & ARI & NMI \\
    \midrule
    Autoencoder     & k-means      & 7  & 0.102  & 0.015 & 0.031 \\
    Autoencoder     & GMM (BIC)    & 16 & $-$0.020 & 0.031 & 0.055 \\
    CNN (scratch)   & k-means      & 7  & 0.306  & 0.008 & 0.016 \\
    CNN (scratch)   & GMM (BIC)    & 12 & $-$0.136 & 0.010 & 0.020 \\
    CNN (warm-start)& k-means      & 7  & 0.044  & 0.016 & 0.024 \\
    CNN (warm-start)& GMM (BIC)    & 16 & $-$0.088 & 0.012 & 0.022 \\
    \bottomrule
  \end{tabular}
\end{table}

\section{Discussion and limitations}
\label{sec:discussion}

Taken together, these results paint a consistent picture: call-type
information exists in all three embedding spaces at a level well above
chance (RQ1), but is not organized in a way that any of the unsupervised
clustering methods we tested can recover (RQ2). We consider this
dissociation -- rather than either result in isolation -- to be the main
finding of this study. It shows that a decodable linear direction for
call-type is not sufficient for an unsupervised clustering algorithm to
discover it unprompted; the relevant information may be encoded in a
low-variance or non-axis-aligned subspace that reconstruction- or
detection-oriented objectives have no incentive to make geometrically
salient. Concretely, this means the unsupervised subdivision of the
"complex call" catch-all category envisioned in prior work
cannot be accomplished by naively clustering an embedding trained for a
different objective (reconstruction or binary detection); it requires an
embedding explicitly optimized to make call-type structure clusterable,
motivating the contrastive- or metric-learning-based objectives we outline
below.

We also note a second consistent pattern that is scientifically
interesting in its own right: the from-scratch binary-detection embedding
is uniformly worse at call-type discrimination than either the
autoencoder or the warm-started embedding, despite performing marginally
\emph{better} than the warm-started model at the binary detection task
itself (per the companion detector paper). This suggests that optimizing
purely for call/non-call separability actively discards information that
is useful for finer-grained discrimination -- a representation-learning
trade-off that, to our knowledge, has not been explicitly quantified for
this dataset before, and that has a direct practical implication: a
detector optimized purely for binary recall should not be assumed to
produce embeddings suitable for downstream sub-type analysis, even though
it is architecturally identical to one that is.

\textbf{Limitations.} This study has several limitations that qualify how
strongly its conclusions should be read. First, we treat the manually
annotated seven-way Type taxonomy as ground truth for the ARI/NMI
clustering-agreement metrics, but the taxonomy's own documentation
describes at least one of its categories as an explicit catch-all for
signals that resisted clean human categorization; a low ARI could
therefore reflect either an uninformative embedding or a manual taxonomy
that does not itself correspond to a naturally clusterable structure --
our data cannot distinguish between these two explanations, and doing so
is an important direction for future work (e.g., via targeted
expert review of specific discovered clusters). Second, we evaluate only
two clustering families (k-means, Gaussian mixtures) at a single random
seed each; more flexible density-based methods (e.g., HDBSCAN) or
repeated-seed stability analysis might recover structure that these two
methods miss, and we plan to report this in a full version of this study.
Third, our linear probe is a single, comparatively simple classifier
(multinomial logistic regression); a nonlinear probe could reveal
additional decodable structure not visible to a linear decision boundary,
which would strengthen rather than weaken the RQ1/RQ2 dissociation
reported here. Fourth, none of the three embeddings evaluated here were
trained with any objective that specifically targets sub-type structure
(e.g., contrastive learning with call-preserving augmentations, or
metric learning against partial sub-type supervision); this is precisely
the gap our findings motivate addressing next.

\textbf{Future work.} The most direct next step is to replace the purely
reconstructive or purely binary-discriminative pretraining objectives
studied here with a representation-learning objective explicitly designed
to make call-type structure separable -- for instance, a contrastive
objective using label-preserving spectrogram augmentations (time-shift,
additive noise at controlled SNR, frequency-axis jitter) to pull same-call
views together and push different-call views apart, or a semi-supervised
metric-learning approach that uses a small fraction of Type-labeled calls
to shape the embedding geometry before applying unsupervised clustering to
the remainder. Any such method should be validated not only by
ARI/NMI against the existing manual taxonomy, but by a human-expert review
of discovered clusters that do \emph{not} match the manual taxonomy, since
-- as noted above -- the manual taxonomy's own "complex call" category was
defined as a catch-all rather than a settled ground truth, and a genuinely
useful method might legitimately propose a different partition than the
one used here for evaluation. Beyond call-type structure, the original
research program motivating this dataset also identified single-spectrogram
range estimation and multi-sensor context fusion (combining detections
across co-located DASAR channels, and explicitly modeling seismic airgun
temporal periodicity to reduce false alarms) as further open objectives;
neither is addressed by the present study or its companion detector paper,
and both remain open targets for future work building on the
representations studied here.

\section*{Broader impact}
Automated tools of this kind are intended to help researchers process
existing, currently under-analyzed decades-scale passive acoustic
archives more efficiently, in support of monitoring bowhead whale
distribution and behavior amid climate change and increasing human
activity in the Arctic -- a topic of direct importance to Alaska Native
communities who rely on the bowhead subsistence hunt. We are not aware of
a plausible negative-use pathway for this specific probing/clustering
study, which does not itself produce a deployable detector or classifier.

\begin{ack}
This research was supported by the North Pacific Research Board. We thank
Greeneridge Sciences, Inc.\ for the manual call annotations underlying
this dataset.
\end{ack}

\begin{thebibliography}{10}

\bibitem{thode2012automated}
A.~M. Thode et al.
\newblock An automated approach for identifying and classifying bowhead whale calls (working citation; verify full bibliographic details before submission -- shared with the companion JASA manuscript).

\bibitem{shiu2020deep}
Y.~Shiu, K.~J. Palmer, M.~A. Roch, E.~Fleishman, X.~Liu, E.-M. Nosal,
T.~Helble, D.~Cholewiak, D.~Gillespie, and H.~Klinck.
\newblock Deep neural networks for automated detection of marine mammal
species.
\newblock \emph{Scientific Reports}, 10:607, 2020.

\bibitem{allen2021convolutional}
A.~N. Allen, M.~Harvey, L.~Rideout, D.~Nason, S.~Perry, and A.~D. Jarvis.
\newblock A convolutional neural network for automated detection of
humpback whale song in a diverse, long-term passive acoustic dataset.
\newblock \emph{Frontiers in Marine Science}, 8, 2021.

\bibitem{ozanich2021deep}
E.~Ozanich, A.~Thode, P.~Gerstoft, L.~A. Freeman, and S.~Freeman.
\newblock Deep embedded clustering of coral reef bioacoustics.
\newblock \emph{Journal of the Acoustical Society of America},
149(4):2587--2601, 2021.

\bibitem{hinton2006reducing}
G.~E. Hinton and R.~R. Salakhutdinov.
\newblock Reducing the dimensionality of data with neural networks.
\newblock \emph{Science}, 313(5786):504--507, 2006.

\bibitem{paszke2019pytorch}
A.~Paszke, S.~Gross, F.~Massa, A.~Lerer, J.~Bradbury, G.~Chanan,
T.~Killeen, Z.~Lin, N.~Gimelshein, L.~Antiga, et al.
\newblock PyTorch: An imperative style, high-performance deep learning
library.
\newblock In \emph{Advances in Neural Information Processing Systems 32},
2019.

\bibitem{kingma2015adam}
D.~P. Kingma and J.~Ba.
\newblock Adam: A method for stochastic optimization.
\newblock In \emph{International Conference on Learning Representations},
2015.

\bibitem{ioffe2015batch}
S.~Ioffe and C.~Szegedy.
\newblock Batch normalization: Accelerating deep network training by
reducing internal covariate shift.
\newblock In \emph{International Conference on Machine Learning}, 2015.

\bibitem{nair2010rectified}
V.~Nair and G.~E. Hinton.
\newblock Rectified linear units improve restricted Boltzmann machines.
\newblock In \emph{International Conference on Machine Learning}, 2010.

\bibitem{lecun1998gradient}
Y.~LeCun, L.~Bottou, Y.~Bengio, and P.~Haffner.
\newblock Gradient-based learning applied to document recognition.
\newblock \emph{Proceedings of the IEEE}, 86(11):2278--2324, 1998.

\bibitem{mcinnes2018umap}
L.~McInnes, J.~Healy, and J.~Melville.
\newblock UMAP: Uniform manifold approximation and projection for
dimension reduction.
\newblock \emph{arXiv:1802.03426}, 2018.

\bibitem{pedregosa2011scikit}
F.~Pedregosa, G.~Varoquaux, A.~Gramfort, V.~Michel, B.~Thirion,
O.~Grisel, M.~Blondel, P.~Prettenhofer, R.~Weiss, V.~Dubourg, et al.
\newblock Scikit-learn: Machine learning in Python.
\newblock \emph{Journal of Machine Learning Research}, 12:2825--2830,
2011.

\bibitem{rousseeuw1987silhouettes}
P.~J. Rousseeuw.
\newblock Silhouettes: A graphical aid to the interpretation and
validation of cluster analysis.
\newblock \emph{Journal of Computational and Applied Mathematics},
20:53--65, 1987.

\bibitem{hubert1985comparing}
L.~Hubert and P.~Arabie.
\newblock Comparing partitions.
\newblock \emph{Journal of Classification}, 2(1):193--218, 1985.

\end{thebibliography}

\end{document}
